% Part: counterfactuals
% Chapter: introduction
% Section: paradoxes-material

\documentclass[../../../include/open-logic-section]{subfiles}

\begin{document}

\olfileid{cnt}{int}{par}

\olsection{বস্তুগত শর্তবচনের আপাত-বৈপরীত্য}

ইংরেজি ``if \dots then \dots'' বক্তব্য প্রতীকায়নে
$!A \lif !B$ ব্যবহারের প্রথম সমালোচকদের অন্যতম ছিলেন
সি.~আই. লুইস। তিনি হোয়াইটহেড ও রাসেলের \emph{Principia
  Mathematica}-তে বস্তুগত শর্তবচনের ব্যবহার নিয়ে আপত্তি
তোলেন; তাঁরা $\lif$-কে ``নিহিত করে'' বলে পড়তেন। লুইস
যথার্থই বলেন, $\lif$-এর অর্থ ``নিহিত করে'' হলে যেকোনো
মিথ্যা বচন~$p$ এমনটি নিহিত করবে যে $p$ থেকে~$q$
নিহিত হয়—কারণ $p$ মিথ্যা হলে $p \lif (p \lif q)$ সত্য।
একইভাবে যেকোনো সত্য বচন~$q$ এমনটি নিহিত করবে যে $p$
থেকে~$q$ নিহিত হয়—কারণ $q$ সত্য হলে
$q \lif (p \lif q)$ সত্য।

অবশ্য যুক্তিবিদেরা জানেন, \emph{নিহিতকরণ}, অর্থাৎ যৌক্তিক
অনুগমন, কোনো সংযোজক নয়; এটি !!{formula}s বা বক্তব্যের
মধ্যকার সম্পর্ক। তাই বিভ্রান্তি এড়াতে $\lif$-কে ``নিহিত
করে'' বলে না পড়াই উচিত।\footnote{গণিতবিদ ও কম্পিউটারবিজ্ঞানীদের
  মধ্যে ``$\lif$''-কে ``নিহিত করে'' বলে পড়ার চল এখনো আছে।
  লুইস যে বিভ্রান্তির কথা তুলেছিলেন, তা এড়াতে দার্শনিকেরা
  একে ``কেবল যদি'' বলে পড়ার চেষ্টা করেন।} তা না পড়লে
লুইসের নির্দিষ্ট আপত্তিটি আর ওঠে না: $p$ দিয়ে ঘটনাচক্রে
মিথ্যা কোনো ইংরেজি বাক্য বোঝালেও $p$, $q$-কে ``নিহিত''
করে না। $p \Entails q$ কি না নির্ধারণে \emph{সব}
!!{valuation}s বিবেচনা করতে হয়। $p$ দিয়ে বাস্তবে মিথ্যা
কোনো বাক্য বোঝালেও $p \Entails/ q$।

তবু লুইসের নিচের বক্তব্যটির মতো ``if \dots then \dots''
বাক্যে অস্বস্তিকর কিছু থেকে যায়:
\begin{quote}
চাঁদ সবুজ পনির দিয়ে তৈরি হলে $2+2=4$।
\end{quote}
একই অস্বস্তি দেখা যায় এই অনুসিদ্ধান্তগুলিতে:
\begin{quote}
  চাঁদ সবুজ পনির দিয়ে তৈরি নয়। অতএব চাঁদ সবুজ পনির
  দিয়ে তৈরি হলে $2+2=4$।

  $2+2 = 4$। অতএব চাঁদ সবুজ পনির দিয়ে তৈরি হলে
  $2+2=4$।
\end{quote}
অথচ ``if \dots then \dots'' নিছক $\lif$ হলে বাক্যটি
নির্বিবাদে সত্য এবং অনুসিদ্ধান্তগুলিও নির্বিবাদে বৈধ হতো।

আরেকটি প্রশ্ন ওঠে সর্বতঃসত্য
$(!A \lif !B) \lor (!B \lif !A)$ থেকে। এর ভিত্তিতে
মনে হতে পারে, সংবাদপত্র থেকে দৈবক্রমে নেওয়া যেকোনো দুটি
নির্দেশক বাক্য~$S$ ও~$T$-এর ক্ষেত্রে ``$S$ হলে $T$,
অথবা $T$ হলে $S$'' সত্য হওয়ার কথা।

\end{document}
